\documentclass[11pt]{article}
\usepackage[T1]{fontenc}
\usepackage[utf8]{inputenc}
\usepackage{lmodern}
\usepackage[a4paper,margin=27mm,headheight=14pt]{geometry}
\usepackage{amsmath,amssymb,amsthm,mathtools}
\usepackage{microtype,booktabs,longtable,array,enumitem,xcolor}
\usepackage{listings,fancyhdr}
\usepackage[colorlinks=true,linkcolor=blue!45!black,citecolor=blue!45!black,urlcolor=blue!45!black]{hyperref}
\hypersetup{pdftitle={Apery's theorem in intuitionistic primitive recursive arithmetic},pdfauthor={Prepared for Vladimir Reshetnikov}}
\pagestyle{fancy}\fancyhf{}\fancyhead[L]{\small Ap\'ery's theorem in iPRA}\fancyhead[R]{\small A finite rational proof}\fancyfoot[C]{\thepage}
\setlength{\parskip}{4pt}\setlength{\parindent}{0pt}
\setlist{nosep,topsep=4pt}
\newtheorem{theorem}{Theorem}[section]
\newtheorem{lemma}[theorem]{Lemma}
\newtheorem{proposition}[theorem]{Proposition}
\newtheorem{corollary}[theorem]{Corollary}
\theoremstyle{definition}\newtheorem{definition}[theorem]{Definition}
\theoremstyle{remark}\newtheorem{remark}[theorem]{Remark}
\newcommand{\N}{\mathbb N}\newcommand{\Z}{\mathbb Z}\newcommand{\Q}{\mathbb Q}
\newcommand{\ipra}{\mathrm{iPRA}}\newcommand{\lc}{\operatorname{lcm}}
\newcommand{\abs}[1]{\left|#1\right|}\newcommand{\floor}[1]{\left\lfloor#1\right\rfloor}
\newcommand{\HH}[2]{H_{#1}^{(#2)}}
\newcommand{\code}[1]{{\small\texttt{#1}}}
\lstset{basicstyle=\ttfamily\small,columns=fullflexible,breaklines=true,frame=single,rulecolor=\color{black!20},showstringspaces=false}
\title{\textbf{Ap\'ery's theorem in intuitionistic\break primitive recursive arithmetic}\\[5pt]
\large A fully finite rational separation argument\\and a Lean formalization blueprint}
\author{Prepared for Vladimir Reshetnikov}
\date{15 November 2026}

\begin{document}
\maketitle
\begin{abstract}
We give a constructive arithmetic proof that $\zeta(3)$ is irrational.
Its actual conclusion is stronger and contains no real numbers or infinite sums:
for every integer $p$ and positive integer $q$, a primitive-recursively specified
finite rational interval containing the partial sums of $\sum j^{-3}$ separates
from $p/q$ by an explicitly specified positive amount.
The proof combines the classical Nesterenko--Zudilin rational function with
an algebraic partial-fraction functional, an exact Wronskian, and the elementary
sparse-index bound $\lc(1,\ldots,42^k)^3\le28^{42^k}$.
All estimates are rational. Every apparent limiting step is replaced by a finite
calculation with explicit auxiliary indices and a decidable rational comparison.
We identify the quantifier-free inductions needed in iPRA and give executable
Lean definitions and a module-by-module plan for constructing an object-language
derivation in ProveIt's ConciseIPRA calculus.
The mathematical argument is developed here; the accompanying Lean files are
uncompiled specifications, not a completed formal proof.
\end{abstract}

\section*{Deliverable and verification status}
\addcontentsline{toc}{section}{Deliverable and verification status}
\begin{center}
\begin{tabular}{p{.26\linewidth}p{.65\linewidth}}
\toprule
Component & Status and exact scope\\\midrule
Mathematical proof & Complete finite rational argument, including the recurrence,
denominator bound, positivity, separation, and induction audit.\\
Exact computation & The supplied Python verifier checks the fixed polynomial
certificate symbolically and checks representative finite identities with
integer and rational arithmetic. These checks supplement the proof.\\
Lean artifacts & Definitions, finite theorem specification, an interface to the
actual ConciseIPRA object theory, and an explicit implementation roadmap.
No claim that the new files have been compiled or that the final internal
derivation has been constructed.\\
Repository dependency & Interfaces inspected at ProveIt commit
\texttt{d72c8b24637e}.
No result in an unverified draft is used as a proved theorem.\\
\bottomrule
\end{tabular}
\end{center}
\newpage
\tableofcontents
\newpage

\section{The exact arithmetic theorem}
Write
\begin{equation}\label{eq:partials}
 Z_M=\sum_{j=1}^{M}\frac1{j^3},\qquad
 a_n=\sum_{k=0}^{n}\binom nk^2\binom{n+k}{k}^2.
\end{equation}
Empty sums have value zero. For $q\ge1$, define the following natural numbers:
\begin{equation}\label{eq:modulus}
 n(q)=42^{q+1},\qquad
 h(q)=(n(q)+1)^3a_{n(q)+1},\qquad
 N(q)=a_{n(q)}h(q).
\end{equation}
Each definition uses only primitive recursion, finite sums, multiplication,
and integer powers. In particular $N(q)$ is a specified function, not the
outcome of an unbounded search.

\begin{theorem}[Finite Ap\'ery separation]\label{thm:main}
For every $p\in\Z$ and $q\in\N$ with $q\ge1$, put $N=N(q)$. Then
\begin{equation}\label{eq:main}
 \boxed{\quad
 \frac pq+\frac1N<Z_N
 \quad\lor\quad
 Z_N+\frac1{2N^2}+\frac1N<\frac pq.
 \quad}
\end{equation}
After signed integers and rational inequalities are coded by natural numbers,
this is a theorem of intuitionistic primitive recursive arithmetic.
\end{theorem}

The interval $[Z_N,Z_N+1/(2N^2)]$ contains $Z_M$ for every $M\ge N$;
this is proved by a finite telescoping inequality in Section~\ref{sec:finite}.
Thus the theorem gives a direction and a positive rational margin separating
$p/q$ from the entire interval. In any usual interpretation of the rational
sequence as the real number $\zeta(3)$, it gives
\begin{equation}\label{eq:apartness}
 \abs{\zeta(3)-\frac pq}>\frac1{N(q)}.
\end{equation}
The huge modulus is chosen for a short proof, not for numerical efficiency or
an optimized irrationality measure. The finite theorem, rather than the real
notation in \eqref{eq:apartness}, is the internal iPRA statement.

\subsection{Proof strategy and attribution}
We construct $b_n\in\Q$ such that $d_n^3b_n\in\Z$, where
$d_n=\lc(1,\ldots,n)$ and $d_0=1$. An exact recurrence gives a positive
Wronskian. A coarse finite estimate identifies the resulting approximants
with the sequence $Z_M$. This yields
\[
 \frac6{(n+1)^3a_{n+1}}-\frac{a_n}{2M^2}
 \le a_nZ_M-b_n\le\frac1{4a_n}\qquad(n\ge48).
\]
The separate bound $d_{42^k}^3\le28^{42^k}$ lets us choose $n=n(q)$ with
$qd_n^3/a_n<1/8$. A decidable comparison of one integer with zero then proves
the disjunction.

The rational function and its telescoper are classical and are recorded in
Zudilin's elementary proof \cite{Zudilin}, following Nesterenko's approach.
The Egyptian-fraction and multinomial idea behind the denominator estimate
is due to Hanson \cite{Hanson}. We prove the precise weaker sparse-index bound
needed here directly. Our presentation supplies the finite formal-functional
recurrence argument, the explicit rational squeeze, and the iPRA separation
certificate. It does not claim discovery of Ap\'ery's theorem or of its
classical approximation sequences.

\section{Foundations and allowed reasoning}\label{sec:foundations}
We use the standard first-order formulation of iPRA. Its variables range over
natural numbers. There is a function symbol for each primitive recursive
function, with its defining equations for zero, successor, projections,
composition and primitive recursion. There are the usual equality rules,
successor injectivity, and $Sx\ne0$. Logic is intuitionistic first-order logic.
Induction is restricted to quantifier-free formulas, with parameters:
\[
 A(0),\qquad A(x)\longrightarrow A(x+1)
 \quad\Longrightarrow\quad A(t).
\]
These ingredients are precisely the object-theory interface used by the
ConciseIPRA files discussed in Section~\ref{sec:lean}.

Integers may be represented by pairs $(u,v)$ for $u-v$. Rationals are pairs
of a signed numerator and a positive denominator. Equality and order are
defined by cross multiplication; reduction to lowest terms is unnecessary.
After moving denominator signs into numerators, rational expressions in
this paper have explicitly positive denominators at the point where order
comparisons are made. Thus every rational inequality
can be replaced by a quantifier-free relation on natural numbers.

Finite sums and products are primitive recursive. So are factorials,
binomial coefficients, natural division, remainders, bounded primality tests,
gcd, and lcm. For example, gcd can be found by a scan of $0,\ldots,u+v$;
positive lcm can then be computed by a bounded quotient of $uv$ by the gcd.
No complexity restriction on running time is imposed by iPRA.

Bounded universal claims can be replaced by primitive recursive Boolean
folds. For instance, ``$P(i)$ for all $i\le r$'', when $P$ is a decidable
arithmetic predicate, is the equation
$\prod_{i=0}^{r}\chi_P(i)=1$. Its usual induction proof becomes an induction
on this quantifier-free equation. This convention also permits finite
arrays, coefficient lists and bounded course-of-values arguments. It does
not license induction on an arbitrary formula containing unbounded
quantifiers.

Every case split below is on natural, integer or rational order, equality,
or a bounded test. These predicates have primitive recursive decision
functions. No instance of excluded middle for an arbitrary proposition,
Markov's principle, countable choice, real completeness, or an unbounded
least-number principle is used.

\section{Finite coefficients and their denominators}\label{sec:coeff}
For $0\le k\le n$ and $s\in\{2,3\}$, put
\[
 c_{n,k}=(-1)^{n-k}\binom nk\binom{n+k}{k},\qquad
 \HH{k}{s}=\sum_{j=1}^{k}\frac1{j^s}.
\]
The signed coefficients $c_{n,k}$ are integers. Define
\begin{align}
 a_n&=\sum_{k=0}^{n}c_{n,k}^2,\label{eq:a}\\
 b_n&=\sum_{k=0}^{n}c_{n,k}^2\HH{k}{3}
 +\sum_{0\le k<l\le n}\frac{c_{n,k}c_{n,l}}{l-k}
       \bigl(\HH{k}{2}-\HH{l}{2}\bigr).\label{eq:b}
\end{align}
Formula \eqref{eq:a} agrees with \eqref{eq:partials}. These definitions,
including all denominators, are finite.

\begin{lemma}[Integrality and initial values]\label{lem:integrality}
For every $n\ge0$, $a_n$ is a positive integer and $d_n^3b_n\in\Z$.
Moreover
\[
 a_0=1,\quad a_1=5,\quad b_0=0,\quad b_1=6.
\]
\end{lemma}
\begin{proof}
The $k=0$ summand of $a_n$ is $1$. For the first sum in $b_n$, every
denominator is $j^3$ with $1\le j\le n$, so it divides $d_n^3$.
For the second sum, each expanded denominator is $(l-k)j^2$ with
$1\le l-k\le n$ and $1\le j\le n$. Both $l-k$ and $j$ divide $d_n$;
explicitly
\[
 \frac{d_n^3}{(l-k)j^2}=\frac{d_n}{l-k}\left(\frac{d_n}{j}\right)^2\in\N.
\]
Thus integrality holds term by term, before any cancellation.
At $n=0$ the harmonic sums vanish. At $n=1$ the coefficients are
$c_{1,0}=-1$, $c_{1,1}=2$, giving $a_1=1+4=5$ and
$b_1=4+(-2)(0-1)=6$.
\end{proof}

The denominator proof can alternatively define an integer $B_n=d_n^3b_n$
directly, replacing each $1/j^3$ by $(d_n/j)^3$ and each
$1/((l-k)j^2)$ by $(d_n/(l-k))(d_n/j)^2$. This is useful when encoding the
final integer comparison without a rational-integrality predicate.

\subsection{The rational function and partial fractions}
Use the formal variable $t$ and define
\begin{equation}\label{eq:R}
 Q_n(t)=\frac{\prod_{j=1}^{n}(t-j)}{\prod_{j=0}^{n}(t+j)},\qquad
 R_n(t)=Q_n(t)^2.
\end{equation}
These are finite polynomial fractions over $\Q$. Formal differentiation
means the usual coefficient rule and quotient rule on these finite objects.
No analytic differentiability theorem is being invoked.

\begin{lemma}[Partial fractions]\label{lem:partial}
We have
\begin{equation}\label{eq:Qpartial}
 Q_n(t)=\sum_{k=0}^{n}\frac{c_{n,k}}{t+k}.
\end{equation}
On putting
\begin{equation}\label{eq:gamma}
 \gamma_{n,k}=\sum_{\substack{0\le l\le n\\l\ne k}}
                  \frac{c_{n,k}c_{n,l}}{l-k},
\end{equation}
we obtain
\begin{equation}\label{eq:Rpartial}
 R_n(t)=\sum_{k=0}^{n}\left(
       \frac{c_{n,k}^2}{(t+k)^2}+\frac{2\gamma_{n,k}}{t+k}\right),
 \qquad\sum_{k=0}^{n}\gamma_{n,k}=0.
\end{equation}
Furthermore
\begin{equation}\label{eq:bgamma}
 b_n=\sum_{k=0}^{n}c_{n,k}^2\HH{k}{3}
                   +\sum_{k=0}^{n}\gamma_{n,k}\HH{k}{2}.
\end{equation}
\end{lemma}
\begin{proof}
Multiply \eqref{eq:Qpartial} by $\prod_{j=0}^{n}(t+j)$. Both sides become
polynomials of degree at most $n$. At $t=-k$, only the $k$th summand on the
right survives. The coefficient it requires is
\[
 \frac{\prod_{j=1}^{n}(-k-j)}{\prod_{0\le j\le n,\,j\ne k}(j-k)}
 =(-1)^{n-k}\frac{(n+k)!}{k!^2(n-k)!}=c_{n,k}.
\]
The difference therefore has $n+1$ distinct roots and degree at most $n$.
Repeated division by the corresponding linear factors proves it zero.
This familiar root-count argument uses finite polynomial division only.

Square \eqref{eq:Qpartial} and, for $k<l$, use the directly verified identity
\[
 \frac1{(t+k)(t+l)}=
 \frac1{l-k}\left(\frac1{t+k}-\frac1{t+l}\right).
\]
This gives \eqref{eq:Rpartial}. The two terms indexed by $(k,l)$ and
$(l,k)$ cancel in $\sum_k\gamma_{n,k}$. Pairing the same two terms in
$\sum_k\gamma_{n,k}\HH{k}{2}$ gives exactly the second sum in
\eqref{eq:b}, proving \eqref{eq:bgamma}.
\end{proof}

\section{An exact recurrence without a limit argument}\label{sec:recurrence}
Let
\begin{equation}\label{eq:telescoper-def}
 P_n=(2n+1)(17n^2+17n+5),\quad
 s_n(t)=4(2n+1)\bigl(-2t^2+t+(2n+1)^2\bigr),\quad
 V_n(t)=s_n(t)R_n(t).
\end{equation}
The following fixed algebraic certificate is the key identity from
\cite{Zudilin}.

\begin{lemma}[Telescoper]\label{lem:telescoper}
For $n\ge1$, as an identity of rational functions,
\begin{equation}\label{eq:telescoper}
 (n+1)^3R_{n+1}(t)-P_nR_n(t)+n^3R_{n-1}(t)
       =V_n(t+1)-V_n(t).
\end{equation}
\end{lemma}
\begin{proof}
The product definitions give the formal quotients
\[
 \frac{R_{n+1}(t)}{R_n(t)}=\left(\frac{t-n-1}{t+n+1}\right)^2,
 \quad
 \frac{R_{n-1}(t)}{R_n(t)}=\left(\frac{t+n}{t-n}\right)^2,
\]
\[
 \frac{R_n(t+1)}{R_n(t)}=
 \frac{t^4}{(t-n)^2(t+n+1)^2}.
\]
Consequently \eqref{eq:telescoper} reduces to the polynomial identity
\begin{align}
 &(n+1)^3(t-n-1)^2(t-n)^2
  -P_n(t-n)^2(t+n+1)^2 \notag\\[-2pt]
 &\hspace{24mm}+n^3(t+n)^2(t+n+1)^2\notag\\
 &\quad=s_n(t+1)t^4-s_n(t)(t-n)^2(t+n+1)^2.\label{eq:polynomial-certificate}
\end{align}
Substituting the two polynomials in \eqref{eq:telescoper-def} and expanding
verifies this identity coefficient by coefficient. It is a single fixed
polynomial identity in $n,t$, not a pattern inferred from sample values.
The supplied verifier expands it exactly. Since we work with rational
functions, intermediate quotients at $t=n$ are never used as numerical
evaluations; the cleared identity establishes the result everywhere its
original denominators are nonzero.
\end{proof}

One must not extract the recurrence for $a_n,b_n$ by assuming that $1$ and
$\zeta(3)$ are linearly independent: that would assume the conclusion.
The following finite device avoids that circularity entirely.

\subsection{A formal partial-fraction functional}
Consider rational functions with constant polynomial part and poles of
order at most two among $0,-1,\ldots,-r$. Write such a function as
\[
 f(t)=c+\sum_{k=0}^{r}\left(\frac{A_k}{(t+k)^2}+\frac{C_k}{t+k}\right).
\]
The decomposition is unique. To see this, clear the common denominator.
Evaluation at each $-k$ first isolates the coefficient of the double pole;
after those coefficients are zero, cancel one copy of each linear factor
and evaluate again to isolate the simple-pole coefficients. The constant
part is then also zero. This proves uniqueness by finite arithmetic.

For two \emph{formal indeterminates} $X,Y$, define
\begin{equation}\label{eq:Theta}
 \Theta(f)=\left(\sum_kA_k\right)X+\frac12\left(\sum_kC_k\right)Y
 -\sum_kA_k\HH{k}{3}-\frac12\sum_kC_k\HH{k}{2}.
\end{equation}
This is just a triple of rational coefficients (of $X,Y,1$), so no polynomial
ring with infinite objects is needed. It is a linear map by uniqueness.
Equations \eqref{eq:Rpartial} and \eqref{eq:bgamma} give
\begin{equation}\label{eq:ThetaR}
 \Theta(R_n)=a_nX-b_n.
\end{equation}

Shifting a pole from $-k$ to $-(k+1)$ changes $\HH{k}{s}$ by
$1/(k+1)^s$. The coefficients of $X,Y$ cancel. Thus direct subtraction gives
\begin{equation}\label{eq:ThetaShift}
 \Theta(f(t+1)-f(t))
 =-\sum_k\frac{A_k}{(k+1)^3}-\frac12\sum_k\frac{C_k}{(k+1)^2}
 =\frac12f'(1).
\end{equation}
Here $f'(1)$ is the evaluation of a formal rational derivative.

The function $V_n=s_nR_n$ is in this class: its polynomial part is constant
because its numerator and denominator have the same degree, and its poles
remain at most double. Its partial fractions can also be constructed
explicitly. Write $s_n(t)=\alpha t^2+\beta t+\delta$ and put
$A_k=c_{n,k}^2$, $C_k=2\gamma_{n,k}$. Multiplication and division of each
partial-fraction term give
\[
 A_k^*=s_n(-k)A_k,\qquad
 C_k^*=s_n'(-k)A_k+s_n(-k)C_k.
\]
The polynomial part is
\[
 \alpha t\sum_kC_k+\alpha\sum_kA_k+
       \sum_kC_k(\beta-\alpha k).
\]
Its potential linear term vanishes because $\sum_kC_k=0$; the remainder is
constant. This explicitly verifies the domain condition for $\Theta$.
For $n\ge1$, $R_n$ has a factor $(t-1)^2$ and a nonzero denominator at $1$,
whence $V_n'(1)=0$.

Apply $\Theta$ to \eqref{eq:telescoper}. By
\eqref{eq:ThetaR}--\eqref{eq:ThetaShift}, the result is a zero formal
polynomial. Compare its $X$ coefficient and constant coefficient.

\begin{proposition}[Ap\'ery recurrence]\label{prop:recurrence}
For either sequence $y=a$ or $y=b$, and every $n\ge1$,
\begin{equation}\label{eq:recurrence}
 (n+1)^3y_{n+1}=P_ny_n-n^3y_{n-1}.
\end{equation}
\end{proposition}
This derivation used finite coefficient identities throughout. In particular,
neither $X$ nor $Y$ denotes a zeta value.

\section{Finite sums identify the approximated number}\label{sec:finite}
We next connect the rational approximants to $Z_M$ without assuming the
existence of an infinite sum.

\begin{lemma}[Finite zeta tails]\label{lem:zeta-tail}
For $T\ge M\ge1$,
\begin{equation}\label{eq:zeta-tail}
 0\le Z_T-Z_M\le\frac1{2M^2}.
\end{equation}
\end{lemma}
\begin{proof}
For $t\ge2$,
\[
 \frac1{t^3}\le\frac12\left(\frac1{(t-1)^2}-\frac1{t^2}\right).
\]
After multiplication by $2t^3(t-1)^2$, the difference of the right numerator
and the left numerator is $3t-2\ge0$. Sum from $t=M+1$ to $T$; the right
side telescopes to $1/(2M^2)-1/(2T^2)$. Nonnegativity is termwise.
\end{proof}

For $n\ge1$ and $T\ge n$, put
\[
 F_{n,T}=-\frac12\sum_{t=1}^{T}R_n'(t),\qquad
 E_n=(n+1)^3\,64^n.
\]

\begin{lemma}[A coarse, uniform finite bound]\label{lem:coarse}
For $n\ge1$ and $T\ge n$,
\begin{equation}\label{eq:coarse}
 \abs{a_nZ_T-b_n}\le3+\frac{E_n}{T^2}.
\end{equation}
\end{lemma}
\begin{proof}
If $1\le t\le n$, the zero of $R_n$ at $t$ is double, so its formal
derivative is zero. For an integer $t>n$, all factors in \eqref{eq:R} are
positive and
\[
 0<R_n(t)\le\frac1{t^2},\qquad
 R_n'(t)=2R_n(t)\left(
       \sum_{j=1}^{n}\frac1{t-j}-\sum_{j=0}^{n}\frac1{t+j}\right).
\]
Every reciprocal inside the parentheses has absolute value at most $1$.
Therefore $\abs{R_n'(t)}\le(4n+2)/t^2$. Also
\[
 \sum_{t=n+1}^{T}\frac1{t^2}
 \le\sum_{t=n+1}^{T}\left(\frac1{t-1}-\frac1t\right)\le\frac1n.
\]
It follows that $\abs{F_{n,T}}\le(2n+1)/n\le3$.

For $s=2,3$, set
$\Delta_{T,k}^{(s)}=\HH{T+k}{s}-\HH{T}{s}$.
Differentiate \eqref{eq:Rpartial} and sum from $1$ to $T$.
All sums are finite. The terms in $\HH{T}{2}$ cancel because
$\sum_k\gamma_{n,k}=0$, and we obtain the exact equality
\begin{equation}\label{eq:finite-pf}
 F_{n,T}=a_nZ_T-b_n+\varepsilon_{n,T},\quad
 \varepsilon_{n,T}=\sum_kc_{n,k}^2\Delta_{T,k}^{(3)}
                       +\sum_k\gamma_{n,k}\Delta_{T,k}^{(2)}.
\end{equation}
Since $0\le\Delta_{T,k}^{(s)}\le k/T^s$, we can bound this error explicitly.
The binomial theorem gives
\[
 \binom nk\le2^n,\qquad \binom{n+k}{k}\le2^{n+k}\le4^n,
\]
so $\abs{c_{n,k}}\le8^n$ and $c_{n,k}^2\le64^n$.
In \eqref{eq:gamma} each nonzero integer $l-k$ has absolute value at least
$1$, so $\abs{\gamma_{n,k}}\le n64^n$. Consequently, for $T\ge1$,
\[
 \abs{\varepsilon_{n,T}}
 \le\frac{n(n+1)64^n}{T^3}+\frac{n^2(n+1)64^n}{T^2}
 \le\frac{(n+1)^3\,64^n}{T^2}.
\]
Combine this with $\abs{F_{n,T}}\le3$.
\end{proof}

Only this very coarse estimate is needed to identify the number being
approximated. The much smaller error used for irrationality will come from
the recurrence, not from estimating the rational function sharply.

\section{Wronskian, growth, and finite positive tails}\label{sec:wronskian}
Write $r_n=b_n/a_n$, which is defined because $a_n\ge1$.

\begin{lemma}[Exact Wronskian]\label{lem:wronskian}
For every $n\ge0$,
\begin{equation}\label{eq:wronskian}
 a_nb_{n+1}-a_{n+1}b_n=\frac6{(n+1)^3}.
\end{equation}
Hence, for $m>n$,
\begin{equation}\label{eq:r-tail}
 r_m-r_n=6\sum_{k=n}^{m-1}\frac1{(k+1)^3a_ka_{k+1}}.
\end{equation}
\end{lemma}
\begin{proof}
Let $W_n$ be the left side of \eqref{eq:wronskian}. The initial values give
$W_0=6$. Substitution of \eqref{eq:recurrence}, once for $a$ and once for
$b$, gives
\[
 (n+1)^3W_n=n^3(a_{n-1}b_n-a_nb_{n-1})=n^3W_{n-1}.
\]
Induction proves $(n+1)^3W_n=6$. Divide by positive $a_na_{n+1}$ and
telescope from $n$ to $m-1$.
\end{proof}

\begin{lemma}[Two lower bounds for $a_n$]\label{lem:growth}
The sequence $a_n$ is nondecreasing and satisfies
\begin{equation}\label{eq:grow4}
 a_n\ge4^n\ge n+1.
\end{equation}
For $n\ge48$,
\begin{equation}\label{eq:grow32}
 a_{n+1}\ge32a_n,\qquad a_n\ge32^{n-48}.
\end{equation}
\end{lemma}
\begin{proof}
For fixed $k\le n$, both nonnegative binomial factors in the $k$th summand
of $a_{n+1}$ are at least the corresponding factors in $a_n$, by Pascal's
identity. The extra summand is nonnegative. Thus $a_{n+1}\ge a_n$.
The $k=n$ summand gives
\[
 a_n\ge\binom{2n}{n}^2,
 \qquad\binom{2n}{n}=\prod_{j=1}^{n}\frac{n+j}{j}\ge2^n.
\]
Here $n+j\ge2j$ for $1\le j\le n$; the case $n=0$ is the empty product.
The elementary inequality $4^n\ge n+1$ follows by induction.

Since $a_{n-1}\le a_n$, the recurrence gives
\[
 (n+1)^3a_{n+1}\ge(P_n-n^3)a_n.
\]
The excess over $32(n+1)^3a_n$ is
\[
 \bigl(n^3-45n^2-69n-27\bigr)a_n.
\]
For $n=48+u$ its polynomial factor is
\begin{equation}\label{eq:growth-poly}
 u^3+99u^2+2523u+3573\ge0.
\end{equation}
This proves the first inequality in \eqref{eq:grow32}; iteration from $48$
and $a_{48}\ge1$ gives the second.
\end{proof}

\begin{lemma}[Bounds on every finite approximant tail]\label{lem:finite-tail}
For $m>n\ge48$, put $h_n=(n+1)^3a_{n+1}$. Then
\begin{equation}\label{eq:finite-tail-bounds}
 \frac6{a_nh_n}\le r_m-r_n\le\frac1{4a_n^2}.
\end{equation}
\end{lemma}
\begin{proof}
The first term of the positive sum \eqref{eq:r-tail} gives the lower bound.
For $k=n+j$, \eqref{eq:grow32} gives
$a_k\ge32^ja_n$ and $a_{k+1}\ge32^{j+1}a_n$. Thus
\[
 r_m-r_n
 \le\frac6{32a_n^2}\sum_{j=0}^{m-n-1}\frac1{1024^j}
 \le\frac{192}{1023a_n^2}<\frac1{4a_n^2}.
\]
The finite geometric-sum bound follows by multiplying by $1-1/1024$ and
telescoping. The final numerical comparison is $768<1023$.
\end{proof}

\section{Eliminating every apparent limiting step}\label{sec:squeeze}
The next lemma is about rational order only.

\begin{lemma}[Finite rational squeeze]\label{lem:squeeze}
Let $x,y\in\Q$. If for every positive integer $K$ one has
$x\le y+1/K$, then $x\le y$. If for every such $K$ one has
$x\ge y-1/K$, then $x\ge y$.
\end{lemma}
\begin{proof}
Rational order is decidable. If $x>y$, write their difference using the
product $v$ of their positive denominators. Its positive integer numerator
is at least $1$, so $x-y\ge1/v$. Substitute the explicit integer $K=2v$ in
the hypothesized inequality to get $x-y\le1/(2v)$, a contradiction.
The lower version follows by interchanging $x,y$.
No convergence theorem or unbounded choice is present: a single computable
instance of the displayed family of inequalities suffices.
\end{proof}

\begin{proposition}[The finite enclosure]\label{prop:enclosure}
For all $n\ge48$ and $M\ge1$, with $h_n=(n+1)^3a_{n+1}$,
\begin{equation}\label{eq:enclosure}
 \boxed{\quad
 \frac6{h_n}-\frac{a_n}{2M^2}
 \le a_nZ_M-b_n\le\frac1{4a_n}.
 \quad}
\end{equation}
\end{proposition}
\begin{proof}
Fix an arbitrary positive integer $K$, and choose explicitly
\begin{equation}\label{eq:aux-cutoffs}
 m=n+1+6K,\qquad
 E_m=(m+1)^3\,64^m,\qquad
 T=M+m+2K(E_m+1).
\end{equation}
These are primitive recursive expressions in $n,M,K$.
They ensure $m>n$, $T\ge M$, and $T\ge m$.
By \eqref{eq:coarse},
\[
 \abs{Z_T-r_m}\le\frac3{a_m}+\frac{E_m}{a_mT^2}.
\]
The first summand is less than $1/(2K)$, since
$a_m\ge m+1>6K$. The second is also less than $1/(2K)$:
$a_m\ge1$ and $T^2\ge T\ge2K(E_m+1)>2KE_m$.
Consequently
\begin{equation}\label{eq:aux-close}
 \abs{Z_T-r_m}<\frac1K.
\end{equation}
Combine this with $Z_M\le Z_T\le Z_M+1/(2M^2)$ and
\eqref{eq:finite-tail-bounds}. It follows that
\[
 r_n+\frac6{a_nh_n}-\frac1{2M^2}-\frac1K
 \le Z_M
 \le r_n+\frac1{4a_n^2}+\frac1K.
\]
All quantities here are rational. Apply Lemma~\ref{lem:squeeze} to each
side, then multiply by the positive integer $a_n$. This proves
\eqref{eq:enclosure}.
\end{proof}

This proposition is the main constructive repair to the usual presentation
with infinite series. The index $m$ is not selected by an unexplained
``sufficiently large'' assertion; it is given in \eqref{eq:aux-cutoffs}.
The same holds for $T$. The parameter $K$ is finally instantiated at twice
an explicitly known denominator in the rational squeeze. The proof is
therefore a finite uniform arithmetic argument even when written with
universally quantified lemmas.

\section{An elementary denominator bound at sparse indices}\label{sec:lcm}
The familiar all-index estimate $d_n<3^n$ is more than we need. The following
short variant of Hanson's multinomial method \cite{Hanson} suffices.

\begin{lemma}[A floor certificate]\label{lem:floor}
For a natural number $j$, let
\[
 w(j)=j-\floor{j/2}-\floor{j/3}-\floor{j/7}-\floor{j/42}.
\]
Then $w(j)\ge0$ for all $j$, and $w(j)\ge1$ when $1\le j\le41$.
\end{lemma}
\begin{proof}
The identity $1/2+1/3+1/7+1/42=1$ shows that the sum of the four floors is
at most $j$. If $1\le j\le41$, the last floor is zero and the other three
sum to at most $41j/42<j$. An integer strictly below $j$ is at most $j-1$.
Multiplication by $42$ makes both arguments integer inequalities.
\end{proof}

For $m\ge1$, define the positive integer
\begin{equation}\label{eq:multinomial}
 M_m=\frac{(42m)!}{(21m)!(14m)!(6m)!m!}
     =\binom{42m}{21m}\binom{21m}{14m}\binom{7m}{6m}.
\end{equation}

\begin{lemma}[The divisibility step]\label{lem:lcm-div}
For $m\ge1$, $d_{42m}$ divides $d_mM_m$.
\end{lemma}
\begin{proof}
Fix a prime $p\le42m$. The exponent of $p$ in a positive integer product
is the sum of its exponents in the factors. In a factorial, count once
for every multiple of $p$, once more for every multiple of $p^2$, and so
on. This gives the finite Legendre formula
\[
 v_p(r!)=\sum_{e=1}^{r}\floor{r/p^e}.
\]
Terms with $p^e>r$ vanish. The exponent of $p$ in $d_r$ is the number of
positive exponents $e\le r$ with $p^e\le r$: the multiples of $p^e$ occur
among $1,\ldots,r$ exactly under that condition. Applying the factorial
formula to \eqref{eq:multinomial} gives
\[
 v_p(M_m)=\sum_{e=1}^{42m}
          w\!\left(\floor{42m/p^e}\right).
\]
Here we used $\floor{\floor{x}/c}=\floor{x/c}$ for positive integer $c$,
which follows immediately from quotient and remainder.

Every summand is nonnegative. Every prime power satisfying
$m<p^e\le42m$ has $1\le\floor{42m/p^e}\le41$, and hence contributes
at least $1$. Therefore
\[
 v_p(M_m)\ge\#\{e:m<p^e\le42m\}
             =v_p(d_{42m})-v_p(d_m).
\]
The prime exponents of $d_{42m}$ are thus no greater than those of
$d_mM_m$, proving divisibility. All primes and exponents considered are
bounded by $42m$. The elementary finite factorization justification is
expanded in Appendix~\ref{app:arithmetic}.
\end{proof}

\begin{lemma}[A multinomial upper bound]\label{lem:multibound}
Let $A=2^{22}3^{15}7^7$. Then $M_m\le A^m$ and $A^3<28^{41}$.
\end{lemma}
\begin{proof}
In the expansion of $(21+14+6+1)^{42m}=42^{42m}$, the term with respective
exponents $21m,14m,6m,m$ is
\[
 M_m\,21^{21m}14^{14m}6^{6m}.
\]
All terms are nonnegative, so it is at most the whole sum. Factorization of
the integer bases gives the exact identity
\[
 42^{42m}=A^m\,21^{21m}14^{14m}6^{6m}.
\]
Cancel the positive common factor to obtain $M_m\le A^m$.

For the numerical comparison, just two small inequalities suffice:
\[
 3^7=2187<2401=7^4,\qquad
 3^{10}=59049<65536=2^{16}.
\]
Thus $3^{45}=(3^7)^5\,3^{10}<7^{20}2^{16}$, and hence
\[
 A^3=2^{66}3^{45}7^{21}<2^{82}7^{41}=28^{41}.
\]
\end{proof}

\begin{proposition}[Sparse lcm bound]\label{prop:sparse}
For every $k\ge0$,
\begin{equation}\label{eq:sparse}
 \boxed{d_{42^k}^{\,3}\le28^{42^k}.}
\end{equation}
\end{proposition}
\begin{proof}
Lemmas~\ref{lem:lcm-div} and \ref{lem:multibound} give
$d_{42m}\le d_mA^m$. If $d_m^3\le28^m$, then
\[
 d_{42m}^3\le d_m^3A^{3m}\le28^m(28^{41})^m=28^{42m}.
\]
Start at $m=1$, where $d_1^3=1\le28$, and induct on $k$.
The induction formula $d(42^k)^3\le28^{42^k}$ is itself quantifier-free.
\end{proof}

\section{The explicit index and the final integer comparison}\label{sec:finish}
We now prove Theorem~\ref{thm:main}.

\begin{lemma}[The chosen index is large enough]\label{lem:index}
For $q\ge1$ and $n=42^{q+1}$,
\begin{equation}\label{eq:index-small}
 n\ge48,\qquad \frac{qd_n^3}{a_n}<\frac18.
\end{equation}
\end{lemma}
\begin{proof}
Certainly $n\ge42^2=1764$. Put $C=32^{48}=2^{240}$.
Equations \eqref{eq:grow32} and \eqref{eq:sparse} give
\[
 \frac{qd_n^3}{a_n}\le qC\left(\frac78\right)^n.
\]
We bound the right side by a finite block comparison. Induction gives
$42^{q-1}\ge q$, so $n\ge1764q$. For $q\ge1$,
\[
 1764q\ge6(q+243),
\]
because the difference is $1758q-1458\ge300$.
The exact inequality $8^6=262144>235298=2\cdot7^6$ yields
\[
 \left(\frac87\right)^n
 \ge\left(\frac87\right)^{6(q+243)}>2^{q+243}.
\]
Consequently
\[
 qC\left(\frac78\right)^n
 <\frac{q2^{240}}{2^{q+243}}
 =\frac{q}{8\cdot2^q}\le\frac18,
\]
where $q\le2^q$ is another elementary induction.
Every power comparison can be read after clearing positive denominators.
\end{proof}

\begin{proof}[Proof of Theorem~\ref{thm:main}]
Fix $p\in\Z$ and $q\ge1$. Abbreviate
\[
 n=42^{q+1},\quad a=a_n,\quad b=b_n,\quad d=d_n,\quad
 h=(n+1)^3a_{n+1},\quad N=ah.
\]
All these quantities except possibly $b$ are positive integers; $d^3b$ is
an integer. Since $a_{n+1}\ge a_n$, we have $h\ge a\ge1$ and $N\ge a^2$.
Define the two rational endpoints of the corresponding linear-form interval:
\[
 L=aZ_N-b,\qquad U=a\left(Z_N+\frac1{2N^2}\right)-b.
\]
By Proposition~\ref{prop:enclosure},
\begin{equation}\label{eq:LU}
 L\ge\frac6h-\frac{a}{2N^2}
   =\frac6h-\frac1{2ah^2}\ge\frac5h>0,
 \qquad
 U\le\frac1{4a}+\frac1{2ah^2}\le\frac3{4a}.
\end{equation}
In particular, using \eqref{eq:index-small},
\begin{equation}\label{eq:small-U}
 qd^3U<\frac3{32}<\frac14,\qquad
 \frac{qd^3}{h}\le\frac{qd^3}{a}<\frac18.
\end{equation}

The number
\begin{equation}\label{eq:J}
 J=d^3(ap-qb)=d^3ap-q(d^3b)
\end{equation}
is an integer. Decidable integer order gives precisely the alternatives
$J\le0$ or $J\ge1$.

If $J\le0$, then $p/q\le b/a$. Therefore
\[
 Z_N-\frac pq\ge Z_N-\frac ba=\frac La
       \ge\frac5{ah}=\frac5N>\frac1N.
\]
This is the first disjunct of \eqref{eq:main}.

If $J\ge1$, use \eqref{eq:small-U} to obtain
\[
 qd^3\left(U+\frac aN\right)
 =qd^3U+\frac{qd^3}{h}<\frac14+\frac18<1\le J.
\]
Substitute the definitions of $U,J$, cancel the positive factors, and get
\[
 Z_N+\frac1{2N^2}+\frac1N<\frac pq.
\]
This is the second disjunct. The construction and the case distinction use
only finite arithmetic, proving the theorem.
\end{proof}

\section{Why this proves constructive irrationality}\label{sec:irrationality}
The theorem provides more information than a contradiction from an
unexplained assumption that $\zeta(3)$ is rational. Both the precision $N(q)$
and the rational margin $1/N(q)$ are explicit. The sign is decided by the
integer $J$ in \eqref{eq:J}, which can be computed directly using the integer
formula for $d_n^3b_n$ after Lemma~\ref{lem:integrality}.

Here is an entirely first-order arithmetization. Say that a rational $p/q$
is the value of the zeta sequence if
\begin{equation}\label{eq:rational-value}
 \forall M\ge1\quad
 Z_M\le\frac pq\le Z_M+\frac1{2M^2}.
\end{equation}
The displayed intervals are the natural lower-sum and tail enclosures.
In the standard interpretation their common value is $\sum_{j\ge1}j^{-3}$.
Theorem~\ref{thm:main} refutes \eqref{eq:rational-value} using the single
explicit instance $M=N(q)$. Thus iPRA proves the negation of the assertion
that any rational has this value.

One can equally use a regular Cauchy name
\[
 z_s=Z_{2^{s+1}}\qquad(s\in\N).
\]
Lemma~\ref{lem:zeta-tail} gives the rational modulus
$\abs{z_t-z_s}\le2^{-2s-3}$ for $t\ge s$.
For any $M\ge N(q)$, the same branch of \eqref{eq:main} separates $p/q$
from $Z_M$ by more than $1/N(q)$. Choosing, for example, $s=N(q)+1$ makes
$2^{s+1}\ge N(q)$ and $2^{-s}<1/N(q)$. Thus the usual arithmetic equality
condition $\forall s\,\abs{z_s-p/q}\le2^{-s}$ also fails at an explicit
index. There is no need to form a completed field of real numbers inside
the first-order theory.

In a constructive real-number presentation, these statements give
\emph{apartness}: $\zeta(3)$ and $p/q$ have a witnessed positive separation.
The real formulation \eqref{eq:apartness} follows from the same interval
estimates, without changing the arithmetic proof.

\section{Detailed iPRA internalization audit}\label{sec:audit}
It is not enough to observe that a theorem is computable or that its witness
is primitive recursive. Correctness must also be proved with the permitted
induction. We therefore list the actual induction and encoding obligations.

\begin{longtable}{@{}p{.21\linewidth}p{.36\linewidth}p{.35\linewidth}@{}}
\toprule
Part & Data and invariant & Permitted proof form\\\midrule
\endfirsthead
\toprule Part & Data and invariant & Permitted proof form\\\midrule\endhead
Arithmetic & Quotient/remainder, gcd, signed pairs, positive fraction
denominators & Loop invariants are equations and inequalities after clearing
denominators. Bounded tests use Boolean folds.\\
Binomials & Pascal recursion, finite row sums, factorial identities &
Quantifier-free induction on a row index; inner indices are bounded and
compiled into finite arrays.\\
Partial fractions & Lists of polynomial coefficients, roots $-k$, and
explicit $c_{n,k},\gamma_{n,k}$ & Coefficient equality and finite polynomial
division; root-count induction is on finite degree, with a supplied root
list.\\
Telescoper & Fixed polynomial identity \eqref{eq:polynomial-certificate} &
A fixed finite ring calculation. Product-ratio identities use finite product
recursion.\\
Functional $\Theta$ & Finite coefficient triples and harmonic sums &
Linearity, uniqueness and shifts reduce to finite sum equations; no functional
quantification in the final object language.\\
Wronskian & $(n+1)^3(a_nb_{n+1}-a_{n+1}b_n)=6$ &
Quantifier-free induction after clearing the displayed rational denominators.\\
Growth & $a_{48+i}\ge32^i$ and
$a_{n+j}\ge32^ja_n$ with $n\ge48$ & Quantifier-free induction on $i,j$.
The start and step use \eqref{eq:growth-poly}.\\
Finite tails & Explicit harmonic or geometric partial sums &
Quantifier-free induction on the number of terms; telescoping identities and
termwise rational comparison.\\
Squeeze & Denominator of a rational gap, $K=2v$, followed by explicit
$m,T$ & Decidable rational order and universal instantiation. No induction on
a formula containing the universal precision quantifier.\\
Prime exponents & Bounded factor lists and sums of indicators of $p^e\le r$ &
Quantifier-free invariants for bounded factorization, product valuations and
factorial counting.\\
Sparse lcm & $d(42^k)^3\le28^{42^k}$ & One quantifier-free induction on $k$,
using the explicit multinomial inequality.\\
Final theorem & An integer $J$ and a Boolean disjunction of two rational
comparisons & A decidable integer comparison, substitution and positive
cross multiplication. No unbounded search.\\
\bottomrule
\end{longtable}

Universal statements in the prose are lemmas with arbitrary parameters,
not the formulas to which induction is applied. For example, to prove a
bound for all coefficients below a degree, the induction invariant is a
primitive recursive Boolean conjunction over that finite range. For the
squeeze, the family indexed by $K$ is obtained from the explicit expressions
\eqref{eq:aux-cutoffs}; the final proof substitutes a specific $K$ into it.
These distinctions are necessary to keep the proof in iPRA rather than
silently appealing to full Heyting arithmetic.

All finite encodings can be eliminated by their explicit primitive recursive
definitions. This produces an ordinary finite iPRA derivation of the
universally closed Boolean separation equation. The preceding table is a
non-formal proof audit and an implementation guide; a machine-checked
derivation still has to be constructed.

\section{Lean formalization plan and repository integration}\label{sec:lean}
\subsection{Pinned interfaces actually inspected}
The accompanying audit records source inspection of ProveIt at commit
\texttt{d72c8b24637e} \cite{ProveIt}.
The declared toolchain is Lean \texttt{v4.32.0} with Mathlib
\texttt{v4.32.0}.

The relevant interfaces are in namespace \code{LeanProofs.ConciseIPRA}:
\begin{center}
\begin{tabular}{p{.36\linewidth}p{.55\linewidth}}
\toprule
Source module & Interface and use\\\midrule
\code{Syntax.lean} & \code{Fm}, \code{Fm.QF}, connectives, quantifiers,
substitution, and scoped syntax.\\
\code{Deduction.lean} & \code{Prf}, intuitionistic natural deduction,
equality rules, weakening and cut.\\
\code{PR.lean} & \code{PR k}, \code{PTm n}, \code{PRAAx}, defining-equation
lemmas and \code{PRA.qfInd}.\\
\code{BiInterpretation.lean} & \code{pra\_theorem\_iff} and the
translation to the equational presentation.\\
\code{NativeBiInterpretation.lean} & \code{NFm.ofPRA},
\code{pra\_iff\_native}, \code{native\_iff\_pra}.\\
\bottomrule
\end{tabular}
\end{center}
The inspected \code{basel\_ipra.zip} draft supplies a useful convention for
raw signed fractions and an internalization roadmap. Its new theorem files
are explicitly uncompiled specifications, so no Basel theorem or analysis
compiler from that archive is assumed here.

\subsection{Two distinct formalization obligations}
The first layer is a theorem about finite natural and rational calculations
in the Lean metatheory. The second is a derivation \emph{inside the encoded
first-order iPRA calculus}. The first does not automatically supply the
second.

The final object-language target has the form
\begin{lstlisting}
([] : List (Fm PTm 0)) |- [PRAAx] aperyFormula
\end{lstlisting}
where the file uses the repository's actual Unicode turnstile notation.
Here \code{aperyFormula} universally closes a quantifier-free equation saying
that the concrete separation checker returns $0$. Its arguments are
$p_+,p_-,q_0$, representing $p=p_+-p_-$ and $q=q_0+1$; this removes the need
for a nonzero-denominator antecedent.

An implementation that proves only \code{Irrational (Real.zeta 3)} or an
equivalent theorem about \code{Real.tsum} does not meet this target.
Likewise, evaluating a checker on finitely many numerals cannot prove its
universal correctness. No semantic truth-to-provability principle is
available or required.

\subsection{Executable specifications included in the archive}
The files in the \code{lean/} directory have the following roles.
\begin{itemize}
\item \code{AperySpec.lean} supplies total finite definitions for signed raw
fractions, exact order, finite sums, binomial coefficients, $a_n,b_n,d_n$,
$n(q),h(q),N(q)$, $Z_M$, and the separation checker. It states the principal
mathematical obligations with their exact arguments and bounds.
\item \code{FiniteAlgebraSpec.lean} defines finite polynomial and rational-function
operations, the partial-fraction arrays and $\Theta$, the derivative sums,
and the explicit auxiliary indices used in the finite squeeze.
\item \code{AlgebraCertificates.lean} gives ordinary Lean proof scripts for
the fixed polynomial identities and small numerical comparisons.
\item \code{ObjectTheoryTarget.lean} gives the interface from a separately
constructed \code{PR 3} program to the object-language sentence and its
internal derivation target. Transport to the concise native calculus is
stated conditional on an actual derivation in iPRA.
\item \code{FORMALIZATION\_ROADMAP.md} gives the dependency order and the
remaining implementation tasks.
\end{itemize}

These files use explicit parameters for missing derivations. They do not
turn a missing proof into an axiom or conceal it behind a \code{sorry}.
They are not presented as a completed Ap\'ery formalization. No Lean, Lake,
or Elan executable was available in the authoring environment, so the new
Lean source has not been compiled here. This limitation does not affect the
exact Python certificate checks or the human proof, but compilation and
object-proof construction remain concrete required work.

\subsection{A completion path with explicit intermediate targets}
The following order avoids a large analysis development.
\begin{enumerate}
\item Implement and verify raw signed rational arithmetic and bounded loops.
Prove denominator positivity and agreement of comparison with cross
multiplication. Specialize each bounded sum to a first-order primitive
recursive program; do not quantify over arbitrary functions in the object
language.
\item Build the binomial, factorial, gcd/lcm and finite polynomial libraries
with internal defining-equation proofs. Establish the specific coefficient
and valuation invariants in the audit table.
\item Formalize Lemmas~\ref{lem:integrality} and \ref{lem:partial}, then check
the fixed polynomial certificate \eqref{eq:polynomial-certificate}.
An ambient \code{ring} proof is a useful independent check; the internal
version must produce an equality derivation from iPRA's arithmetic lemmas.
\item Represent a partial-fraction expression by finite coefficient arrays,
with \code{Theta} returning a triple of raw fractions. Prove uniqueness,
linearity, and the shift identity for these arrays. Derive both recurrences
by coefficient projection.
\item Prove the finite derivative-sum identity and error estimate, the
Wronskian identity, and the two growth bounds. Formal derivatives here are
operations on polynomial pairs, so no Mathlib real-analysis dependency is
needed.
\item Prove the finite enclosure with the exact indices
\eqref{eq:aux-cutoffs}. Implement rational squeeze as the explicit
denominator-based contradiction, with decidable order supplied internally.
\item Prove the floor certificate, bounded valuation calculation and
multinomial selected-term inequality. Apply the quantifier-free induction
for \eqref{eq:sparse}; prove the small block comparisons of
Lemma~\ref{lem:index}.
\item Construct the integer $J$, perform its internally justified decidable
case split, and derive each Boolean branch of the checker. Universally
introduce the three natural arguments.
\item Compose these derivations into the final \code{Prf PRAAx} object.
Only then use \code{pra\_iff\_native} to transport it. Optionally add a
separate semantic corollary about the usual real $\zeta(3)$.
\end{enumerate}
At every step, the intended result is an internal proof as well as any
convenient ambient Lean theorem. The compilation bridge is an obligation
with explicit correctness proofs, not an automatic feature of the
repository that has been assumed without inspection.

\section{Reproduction and scope of the computational checks}\label{sec:checks}
The archive contains \code{code/verify.py}. Run
\begin{lstlisting}
python3 code/verify.py
\end{lstlisting}
with Python and SymPy installed. The verifier uses exact integers, rational
fractions and symbolic polynomial expansion. It checks the telescoper,
the shifted growth polynomial, the small numerical comparisons, the
definitions and recurrences for initial ranges, denominator integrality,
partial-fraction finite-sum errors, and instances of the lcm construction.
Its output is saved in \code{code/verification\_report.json}.

The fixed symbolic polynomial certificate establishes an identity for its
formal variables. The other finite-range checks are regression checks;
their role is to catch transcription, sign and normalization errors.
Universal validity of the theorem is supplied by the arguments and
inductions above. In particular, the program does not attempt to enumerate
the astronomically long sum $Z_{N(q)}$ from the final theorem.

The PDF is compiled from this \LaTeX{} source with
\begin{lstlisting}
latexmk -pdf -interaction=nonstopmode -halt-on-error apery_ipra.tex
\end{lstlisting}
The source requires standard AMS and typography packages. A separate
README lists the files and their verification status.

\appendix
\section{Elementary arithmetic behind the valuation argument}\label{app:arithmetic}
We give additional detail because an unqualified invocation of a strong
number-theory library would obscure the foundational bound.

\subsection{Bounded factorization}
Primality of $p$ is the bounded condition $p\ge2$ and that no integer
$j$ with $2\le j<p$ divides $p$. For $r>1$, scan $2,\ldots,r$ for its
first divisor. That divisor is prime: a proper factor of it would be a
smaller divisor of $r$. Repeatedly divide by such a factor. Each nontrivial
division reduces the remaining number, so $r$ iterations suffice, with
idle iterations after the remainder reaches $1$. A primitive recursive
state stores the remainder, factor list and accumulated product. The
invariant ``accumulated product times remainder equals the original input''
and the primality tests are quantifier-free after bounded Boolean coding.
This supplies a factorization together with its explicit certificate.

Euclid's lemma follows from bounded Bezout computation: if a prime $p$
does not divide $u$, gcd$(p,u)=1$, and the extended Euclidean algorithm
produces integers $x,y$ with $xp+yu=1$. Multiplication by $v$ proves that
$p\mid uv$ implies $p\mid v$. The Euclidean remainder decreases at each
nonzero step, so a loop bounded by the initial sum of its arguments is
sufficient. Back substitution and the Bezout invariant use signed pairs.
Repeated application proves uniqueness of the exponent of each prime in
a finite product. Thus one positive integer divides another exactly when
each prime exponent of the first is no larger than the corresponding
exponent of the second. For a given comparison it is enough to scan primes
up to the larger input.

\subsection{Finite valuation formulas}
For a prime $p$ and positive integer $r$, define its valuation by the bounded
sum
\[
 v_p(r)=\sum_{e=1}^{r}[p^e\mid r],
\]
where brackets denote the Boolean indicator. Divisibility by powers of a
prime is an initial segment of the exponent list; product valuation follows
from Euclid's lemma and finite factorization. Hence
\[
 v_p(r!)=\sum_{j=1}^{r}v_p(j)
        =\sum_{e=1}^{r}\#\{j\le r:p^e\mid j\}
        =\sum_{e=1}^{r}\floor{r/p^e}.
\]
The middle equality exchanges two finite sums, not infinite series.
For $d_r$, each $j\le r$ divides $d_r$, and the maximum exponent required
is exactly the number of $e\le r$ with $p^e\le r$. This proves the lcm
valuation formula used in Section~\ref{sec:lcm}. All witness searches and
all sums have the displayed natural bounds.

\subsection{Binomial and multinomial facts}
Define binomial coefficients by Pascal recursion, with out-of-range values
zero. Induction gives the finite binomial theorem, nonnegativity,
monotonicity in the upper argument and row sum $2^n$.
Factorial formulas follow by multiplying Pascal's identity by the relevant
factorials and induction. Repeated application of the binomial theorem
gives the four-term multinomial expansion. Its coefficient at
$(21m,14m,6m,m)$ is the product of the three binomial coefficients in
\eqref{eq:multinomial}. The selected-term bound therefore follows from a
finite sum of nonnegative integers. None of these steps requires an
asymptotic estimate for factorials or primes.

\section{A concrete coding of the final sentence}\label{app:coding}
One convenient raw rational representation is a triple
$(u,v,d)$ with natural $u,v$ and $d\ge1$, denoting $(u-v)/d$.
Addition, negation and multiplication are given by
\begin{align*}
 (u,v,d)+(u',v',d')&=(ud'+u'd,vd'+v'd,dd'),\\
 -(u,v,d)&=(v,u,d),\\
 (u,v,d)(u',v',d')&=(uu'+vv',uv'+vu',dd').
\end{align*}
The comparison $(u,v,d)<(u',v',d')$ is the natural inequality
\[
 ud'+v'd<u'd+vd'.
\]
Division by a known positive natural multiplies the denominator. A negative
integer coefficient is represented by swapping the two numerator fields.
These operations preserve positive denominators by direct induction.

For inputs $(p_+,p_-,q_0)$, let $q=q_0+1$, compute the numbers in
\eqref{eq:modulus}, and compute $Z_N$ by a primitive recursive fold. Form
the two strict comparisons in \eqref{eq:main} by the displayed natural
cross-multiplication rule. Return $0$ if their Boolean disjunction is true,
and $1$ otherwise. Call this concrete function $\operatorname{sep}$.
Then an entirely quantifier-free matrix for our theorem is
\[
 \operatorname{sep}(p_+,p_-,q_0)=0.
\]
Universally close its three free variables to obtain the first-order
sentence used as the final formalization target. The existence of the
large intermediate numbers is guaranteed by their primitive recursive
defining equations; no minimization or higher-type choice is involved.

\begin{thebibliography}{9}
\bibitem{Apery}
R. Ap\'ery, \emph{Irrationalit\'e de $\zeta(2)$ et $\zeta(3)$},
Ast\'erisque \textbf{61} (1979), 11--13.
\url{https://www.numdam.org/item/AST_1979__61__11_0/}.
\bibitem{Zudilin}
W. Zudilin, \emph{An elementary proof of Ap\'ery's theorem},
arXiv:math/0202159, 17 February 2002.
The rational function and telescoper used here are equations (1), (5), (6).
\url{https://www.math.ru.nl/~wzudilin/PS/apery.pdf}.
\bibitem{Hanson}
D. Hanson, \emph{On the product of the primes}, Canadian Mathematical
Bulletin \textbf{15}(1) (1972), 33--37.
\href{https://doi.org/10.4153/CMB-1972-007-7}{doi:10.4153/CMB-1972-007-7}.
\bibitem{ProveIt}
V. Reshetnikov, \emph{ProveIt: ConciseIPRA}, repository snapshot
\texttt{d72c8b24637e}, inspected 15 November 2026.
\href{https://github.com/VladimirReshetnikov/ProveIt/tree/d72c8b24637e342816dc92d11bef337748683e51/Logic/PrimitiveRecursiveArithmetic/ConciseIPRA}{Pinned ConciseIPRA source snapshot}. The accompanying repository audit records
the full commit identifier and exact source-file URLs.
\bibitem{Target}
\emph{Ap\'ery's theorem}, Wikipedia, user-specified overview,
consulted 15 November 2026. Used to identify the requested theorem;
the proof above is independently developed from the cited primary work.
\url{https://en.wikipedia.org/wiki/Ap%C3%A9ry%27s_theorem}.
\end{thebibliography}
\end{document}
